% Standalone problem document, generated from the adjacent README.md.
% Compile with XeLaTeX. Mathematical content is maintained in Markdown.
\documentclass[11pt,a4paper]{article}
\usepackage[fontsize=10.5pt]{scrextend}
\usepackage[margin=22mm,headheight=15pt,headsep=9mm,footskip=11mm]{geometry}
\usepackage{fontspec}
\setmainfont{texgyrepagella}[Extension=.otf,UprightFont=*-regular,BoldFont=*-bold,ItalicFont=*-italic,BoldItalicFont=*-bolditalic]
\setsansfont{texgyreheros}[Extension=.otf,UprightFont=*-regular,BoldFont=*-bold,ItalicFont=*-italic,BoldItalicFont=*-bolditalic]
\setmonofont{texgyrecursor}[Extension=.otf,UprightFont=*-regular,BoldFont=*-bold,ItalicFont=*-italic,BoldItalicFont=*-bolditalic,Scale=MatchLowercase]
\usepackage{amsmath,amssymb}
\usepackage{unicode-math}
\setmathfont{texgyrepagella-math.otf}
\usepackage{xcolor,microtype,enumitem,needspace,fancyhdr}
\usepackage{bookmark}
\definecolor{ink}{HTML}{17344B}
\definecolor{accent}{HTML}{197D80}
\definecolor{muted}{HTML}{596773}
\hypersetup{colorlinks=true,linkcolor=accent,urlcolor=accent,pdftitle={AC-02 - Exact
bilinear rank of the 3 × 3 matrix
product},pdfauthor={Open Problems in Numerical Linear Algebra}}
\urlstyle{same}
\setlength{\parindent}{0pt}
\setlength{\parskip}{5pt plus 1pt minus 1pt}
\linespread{1.04}
\setlength{\emergencystretch}{3em}
\widowpenalty=10000
\clubpenalty=10000
\displaywidowpenalty=10000
\allowdisplaybreaks[1]
\setlist{nosep,leftmargin=1.5em}
\providecommand{\tightlist}{\setlength{\itemsep}{0pt}\setlength{\parskip}{0pt}}
\providecommand{\pandocbounded}[1]{#1}
\pagestyle{fancy}
\fancyhf{}
\fancyhead[L]{\sffamily\footnotesize\color{muted}OPEN PROBLEMS IN NUMERICAL LINEAR ALGEBRA}
\fancyhead[R]{\sffamily\bfseries\footnotesize\color{accent}AC-02}
\fancyfoot[L]{\parbox[b]{0.9\textwidth}{\sffamily\fontsize{8}{10}\selectfont\color{muted}Editorial ratings; a bounded search cannot certify that a problem remains unsolved.\\Literature check: 2026-09-08}}
\fancyfoot[R]{\sffamily\footnotesize\color{muted}\thepage}
\renewcommand{\headrulewidth}{0.3pt}
\renewcommand{\headrule}{\hbox to\headwidth{\color{accent}\leaders\hrule height \headrulewidth\hfill}}
\makeatletter
\renewcommand{\section}{\Needspace{5\baselineskip}\@startsection{section}{1}{0pt}{12pt plus 2pt minus 2pt}{4pt}{\sffamily\bfseries\large\color{ink}}}
\renewcommand{\subsection}{\Needspace{4\baselineskip}\@startsection{subsection}{2}{0pt}{9pt plus 1pt minus 1pt}{3pt}{\sffamily\bfseries\normalsize\color{ink}}}
\makeatother
\setcounter{secnumdepth}{0}
\begin{document}
{\sffamily\small\color{muted}Arithmetic and complexity\par}
\vspace{3pt}
{\raggedright\hyphenpenalty=10000\exhyphenpenalty=10000\sffamily\bfseries\fontsize{21}{24}\selectfont\color{ink}Exact
bilinear rank of the \(3\times3\) matrix product\par}
\vspace{7pt}
{\sffamily\small\textbf{Difficulty:} extreme\quad\textbf{Importance:} interesting
to the community\par}
\vspace{4pt}
{\color{accent}\hrule height 0.8pt}
\vspace{6pt}
\textbf{Topic:} small matrix multiplication algorithms

\textbf{Status:} open; admitted after a bounded literature search found
no resolution.

\subsection{Problem statement}\label{problem-statement}

Determine the least integer \(r\) for which there are complex
coefficients \(u_{t,ij},v_{t,ij},w_{ij,t}\) satisfying, for every pair
of complex \(3\times3\) matrices,

\[
(AB)_{ij}=\sum_{t=1}^{r} w_{ij,t}
 \left(\sum_{a,b=1}^{3}u_{t,ab}A_{ab}\right)
 \left(\sum_{c,d=1}^{3}v_{t,cd}B_{cd}\right).
\]

Scalar linear combinations are free in this bilinear-rank model. The
classical upper bound is 23; deciding whether 22 products suffice is
part of determining the exact minimum. Algorithms mixing entries of both
inputs within a factor are outside this specified model.

\subsection{Why it matters}\label{why-it-matters}

Such identities can be applied recursively to matrix blocks.

\subsection{References and status}\label{references-and-status}

Bläser,
\href{https://theoryofcomputing.org/articles/gs005/gs005.pdf}{2013}, §1
and Example 5.10. Y. Sun,
\href{https://arxiv.org/abs/2604.27645}{\emph{An Exact 56-Addition,
Rank-23 Scheme for General 3×3 Matrix Multiplication}} (2026), gives
another 23-product scheme. C. Wang,
\href{https://arxiv.org/abs/2603.07280}{\emph{Automated Lower Bounds for
Bilinear Complexity over Finite Fields}}, v10 (2026), improves a lower
bound over \(\mathbb F_2\); that is not a complex-field resolution.
Searches for ``rank 22 matrix multiplication 2026 proof'' found no
qualifying algorithm or matching lower bound. \textbf{Admitted: no
resolution located.}
\end{document}
